
\medskip
\noindent
\SGARefTarget{sga7.sga7i.eix.n5-3}{5.3}. Using the orthogonality correspondence between strict
pro-Barsotti--Tate subgroups (in the sense specified above) of the dual
pro-groups $T_\ell(A_K)$ and $T_\ell(A'_K)$, one formally deduces
from (\SGARefLink{sga7.sga7i.eix.n5-2-2}{5.2.2}) that the pairing (\SGARefLink{sga7.sga7i.eix.n1-0-2}{1.0.2}) between these pro-groups induces a
pairing
\begin{equation}\tag{\SGARefTarget{sga7.sga7i.eix.n5-3-2}{5.3.2}}
 T_\ell(A_K)^{ab}\times T_\ell(A'_K)^{ab}
 \longrightarrow T_\ell(\mathbb G_{mK})
\end{equation}
which is a \underline{perfect duality}, that is, it identifies each of
the two factors on the left with the dual of the other (in the sense of
\SGARefLink{sga7.sga7i.eix.n5-2-3}{5.2.3}).\footnote{The superscripts $ab$ denote the \emph{abelian
parts}, the quotient of the fixed part by the toric part.}  It is also
clear that (\SGARefLink{sga7.sga7i.eix.n5-3-2}{5.3.2}) is induced by the pairing
\begin{equation}\tag{\SGARefTarget{sga7.sga7i.eix.n5-3-3}{5.3.3}}
 T_\ell(A^o)^{ab}\times T_\ell(A^{\prime o})^{ab}
 \longrightarrow T_\ell(\mathbb G_{mS})
\end{equation}
obtained from the pairing (\SGARefLink{sga7.sga7i.eix.n5-2-4-2}{5.2.4.2}) by passing to the quotients.  Since
this is a perfect duality on the generic fibres, that is, since the
corresponding homomorphism
\begin{equation}\tag{\SGARefTarget{sga7.sga7i.eix.n5-3-4}{5.3.4}}
 T_\ell(A^o)^{ab}
 \longrightarrow D\bigl(T_\ell(A^{\prime o})^{ab}\bigr)
\end{equation}
is an isomorphism on the generic fibres, one naturally wishes to deduce
that it is so over the whole of $S$, or equivalently on the special
fibres.  This deduction is trivially valid if $\ell\neq p$.  It is
also valid, at least in unequal characteristics when $\ell=p$, by a
deep theorem of Tate [33, Th.\ 4].  We shall nevertheless establish it
directly below in \SGARefLink{sga7.sga7i.eix.n7-4-6}{7.4.6}, without using Tate's theorem.

To make precise the geometric meaning of the result we wish to prove,
observe that the pairing obtained from (\SGARefLink{sga7.sga7i.eix.n5-3-3}{5.3.3}) by passage to the
special fibres is precisely the pairing
\begin{equation}\tag{\SGARefTarget{sga7.sga7i.eix.n5-3-5}{5.3.5}}
 T_\ell(B_o)\times T_\ell(B'_o)
 \longrightarrow T_\ell(\mathbb G_{mk})
\end{equation}
deduced from the biextension $W_o$ of
$(A_o^o,A_o^{\prime o})$ by passage to the quotient, where
$B_o$ and $B'_o$ denote respectively the abelian parts of
$A_o^o$ and $A_o^{\prime o}$.  By VIII \SGARefLink{sga7.sga7i.eviii.n4-8}{4.8}, the biextension
$W_o$ comes from a biextension $W_o^{ab}$ of
$(B_o,B'_o)$ by $\mathbb G_{mk}$, unique up to a unique
isomorphism.  Thus the pairing (\SGARefLink{sga7.sga7i.eix.n5-3-5}{5.3.5}) may also be interpreted as the
one deduced from the divisorial correspondence $W_o^{ab}$, and the
corresponding homomorphism
\begin{equation}\tag{\SGARefTarget{sga7.sga7i.eix.n5-3-6}{5.3.6}}
 T_\ell(B_o)
 \longrightarrow D\bigl(T_\ell(B'_o)\bigr)
 \simeq T_\ell\bigl(D(B'_o)\bigr)
\end{equation}
(where, in the last term, $D$ denotes the dual abelian scheme) is
precisely the morphism obtained by applying $T_\ell$ functorially to
the homomorphism
\begin{equation}\tag{\SGARefTarget{sga7.sga7i.eix.n5-3-7}{5.3.7}}
 B_o\longrightarrow D(B'_o)
\end{equation}
defined by $W_o^{ab}$.  Consequently, by the standard dictionary
(cf.\ [5]), saying that (5.3.5) is a perfect duality is also saying that
the kernel of (\SGARefLink{sga7.sga7i.eix.n5-3-7}{5.3.7}), which is a finite group over $k$ by \SGARefLink{sga7.sga7i.eix.n1-5}{1.5}, has
zero $\ell$-primary component.  Saying that this holds for every
$\ell$ therefore means that the isogeny (\SGARefLink{sga7.sga7i.eix.n5-3-7}{5.3.7}) is an isomorphism,
that is, that the divisorial correspondence $W_o^{ab}$ is a perfect
duality between $B_o$ and $B'_o$.  The preceding considerations
show only that the kernel of (\SGARefLink{sga7.sga7i.eix.n5-3-7}{5.3.7}) is a $p$-group.  We therefore
state the more complete result.

\medskip
\noindent
\underline{Theorem} \SGARefTarget{sga7.sga7i.eix.n5-4}{5.4}. \underline{Let $S$ be a trait, and let
$A_K$ be an abelian scheme over $K$ having semistable reduction
over $S$.  Let $B_o$ and $B'_o$ denote respectively the
abelian parts of the Neron models of $A_K$ and of the dual abelian
scheme $A'_K$.  Then the canonical divisorial correspondence
$W_o^{ab}$ on $(B_o,B'_o)$, deduced by VIII \SGARefLink{sga7.sga7i.eviii.n4-8}{4.8} from the
canonical biextension $W_o$ of $(A_o^o,A_o^{\prime o})$ (the
special fibre of the biextension $W$ of \SGARefLink{sga7.sga7i.eix.n1-4}{1.4}), is a perfect duality
between $B_o$ and $B'_o$; that is, it defines an isomorphism
of each of these abelian schemes with the dual of the other.}

The proof of \SGARefLink{sga7.sga7i.eix.n5-4}{5.4} in the general case will be given in section 7, as a
consequence of a very natural and important construction due to
Raynaud.  Let us merely observe here that, in proving \SGARefLink{sga7.sga7i.eix.n5-4}{5.4}, one may
assume that $S$ is complete, by \SGARefLink{sga7.sga7i.eix.n3-3}{3.3}.
